\chapter{Benchmarks methodology (repeatability + fairness)}
\label{ch:benchmarks}

Benchmarks are where performance stories can go wrong fast.
This chapter documents Relax's benchmark harness under \texttt{benchmarks/} and the methodology guidelines captured in \texttt{docs/benchmarks.md}.

The core goals are:

\begin{itemize}
\item \textbf{Repeatability}: seeded randomness, stable profiles, and a clear warmup/measure separation.
\item \textbf{Fairness}: compare like-for-like work (same DOM outcomes, similar update patterns), and control build modes (dev vs prod).
\item \textbf{Interpretability}: report metrics that explain tail latency (p95) and not just averages.
\end{itemize}

\section{Benchmarks run in a real browser}

The harness is designed to run in a browser environment:

\begin{itemize}
\item Uses \texttt{requestAnimationFrame} to align work with frames.
\item Uses \texttt{performance.now()} for timing.
\item Bundles through Rollup configs that target \texttt{browser: true}.
\end{itemize}

See \texttt{benchmarks/index.html} and the entry script \texttt{benchmarks/main.ts}.

\subsection{A small but important detail: Node/jsdom fallback exists}

The harness is designed for browsers first, but \texttt{benchmarks/harness.ts} contains a compatibility fallback:

\begin{itemize}
\item \texttt{now()} prefers \texttt{performance.now()} but falls back to \texttt{Date.now()}.
\item \texttt{nextFrame()} prefers \texttt{requestAnimationFrame} but falls back to \texttt{setTimeout(..., 16)}.
\end{itemize}

This is \emph{not} intended to make Node-based benchmarking ``accurate''. It's there so that small smoke runs and tooling can still function in non-browser environments.

\section{The harness contract}

The core harness is \texttt{benchmarks/harness.ts}.

\subsection{Benchmark case shape}

A benchmark case is a \texttt{BenchmarkCase}:

\begin{itemize}
\item \texttt{name}: stable identifier for reporting
\item \texttt{setup(host)}: returns a controller with\newline
  \texttt{tick()} -> perform one unit of work\newline
  \texttt{teardown()} -> clean up resources and DOM
\item optional \texttt{commit()}: a barrier that resolves only once the framework has committed for the tick
\end{itemize}

The \texttt{commit()} hook is how the harness stays fair across frameworks that batch work differently (e.g. React flush/paint boundaries).

\subsection{Tiny contract for a fair \texttt{BenchmarkCase}}

To keep comparisons honest, the benchmark cases should follow a simple contract:

\begin{itemize}
\item \texttt{tick()} should perform one repeatable unit of work and avoid random extra allocations.
\item \texttt{teardown()} must remove event handlers / timers / subscriptions (avoid cross-case leaks).
\item If a framework batches updates asynchronously, provide \texttt{commit()} that resolves after work is actually committed.
\end{itemize}

\subsection{Warmup vs measure windows}

\texttt{runBenchmark()} splits execution into two windows:

\begin{itemize}
\item \textbf{Warmup frames} (default 30): let JIT, memoization, layout caching, and internal queues settle.
\item \textbf{Measure frames} (default 120): record per-tick durations.
\end{itemize}

Implementation details in \texttt{runBenchmark()}:

\begin{itemize}
\item For each tick: await \texttt{nextFrame()} , record \texttt{t0} , run \texttt{tick()} , await \texttt{commit()} , compute \texttt{dt}.
\item A \textbf{frame drop} is counted when \texttt{dt > frameBudgetMs} (default 16.7ms).
\end{itemize}

\subsection{Reported metrics: avg + p95 + drops}

The harness reports:

\begin{itemize}
\item \texttt{avgMs} , average time per tick
\item \texttt{p95Ms} , 95th percentile tick time (tail latency)
\item \texttt{frameDrops} , number of frames over budget
\end{itemize}

Percentiles are computed on sorted samples with a simple index selection.
This is good enough for stable comparisons at the harness scale.

The harness computes p95 with:\\
	exttt{idx = floor(p * (n - 1))} (clamped).
This is a simple deterministic estimator, which is a feature here: the same sample list always produces the same percentile.

\section{Profiles and seeded runs}

Repeatability relies on standard profiles and seeded randomness.

\subsection{Profiles via URL parameters}

The profile parser lives in \texttt{benchmarks/profile.ts}.

It reads query params from \texttt{benchmarks/index.html}:

\begin{itemize}
\item \texttt{profile}: \texttt{quick | default | stress}
\item \texttt{seed}: integer seed for randomized update patterns
\item \texttt{warmup}, \texttt{frames}, \texttt{budget}: overrides for fine-grained control
\end{itemize}

Defaults are intentionally ``reasonable and stable'' for local runs.

Implementation note: the profile parser clamps overrides to avoid accidental runaway configurations:

\begin{itemize}
\item \texttt{warmup} is clamped to \texttt{[0, 10000]}
\item \texttt{frames} is clamped to \texttt{[1, 100000]}
\end{itemize}

Seed values are forced into uint32 range using \texttt{>>> 0}, which keeps the PRNG behavior consistent.

\subsection{Seeded RNG}

Randomized patterns (e.g. which list indices change each tick) should be deterministic.
Benchmark code uses a seeded RNG in \texttt{benchmarks/rng.ts}.

A key methodology rule: the seed must be part of the report output so results can be reproduced.

The benchmark RNG implementation is \texttt{Mulberry32} (\texttt{benchmarks/rng.ts}). It is simple, fast, and deterministic---perfect for selecting indices or toggling boolean flags in a way that remains repeatable across runs.

\section{The benchmark runner UI}

The main entrypoint \texttt{benchmarks/main.ts} wires a button click to sequential benchmark runs.

Key details:

\begin{itemize}
\item It resets \texttt{host.innerHTML = ''} between runs so each case starts clean.
\item It prints formatted results to a \texttt{<pre>}.
\item It runs both HRBR and VDOM variants, and includes adapters for React/Solid.
\end{itemize}

Because different frameworks have different setup/teardown costs, the harness keeps the per-tick loop stable and encourages cases to isolate the steady-state behavior.

\subsection{Why \texttt{host.innerHTML = ''} between runs matters}

Clearing between runs gives you two strong properties:

\begin{itemize}
\item \textbf{Isolation:} DOM nodes and event listeners from a previous framework/case do not bleed into the next.
\item \textbf{Repeatability:} each case starts from the same initial DOM shape.
\end{itemize}

This is also why safe \texttt{teardown()} behavior is part of the case contract: the harness clears DOM, but frameworks may have timers/schedulers that need explicit disposal.

\section{Cases: what they measure}

Cases live under \texttt{benchmarks/cases/}.
The intent is to cover:

\begin{itemize}
\item \textbf{Pure text throughput}: \texttt{text-1m}
\item \textbf{Large lists with sparse updates}: \texttt{list-10k-1pct}
\item \textbf{Many small components/widgets}: \texttt{widgets-200}
\item Extra comparisons in \texttt{cases/more.ts}
\end{itemize}

Important fairness guideline: each case should produce the same DOM output across frameworks (or as close as is reasonable).

\subsection{Case study: \texttt{text-1m}}

In \texttt{benchmarks/cases/text-1m.ts}, the HRBR variant builds the smallest possible compiled block:

\begin{itemize}
\item \texttt{templateHTML: `<div><span> </span></div>`} includes an explicit space so there is a real text node.
\item Slot path \texttt{[0,0]} means: root \texttt{<div>} child 0 is \texttt{<span>}, and its child 0 is the \texttt{Text} node.
\item \texttt{tick()} updates a signal in chunks of \texttt{10,000} so work remains frame-visible.
\end{itemize}

This isn't a ``real app'' workload, but it is an excellent microbenchmark for measuring pure slot-write throughput.

\subsection{Case study: \texttt{list-10k-1pct}}

In \texttt{benchmarks/cases/list-10k-1pct.ts}, the cases delegate to example mount helpers:

\begin{itemize}
\item HRBR: \texttt{mountList10k1pctHRBR(host, 10000)}
\item VDOM: \texttt{mountList10k1pctVDOM(host, 10000)}
\end{itemize}

The contract is the same for both: return a controller with \texttt{update1pct()} and \texttt{dispose()}.
Methodologically, that matters: both implementations are being asked to do the same "one tick" of work.

\section{Build modes matter (dev vs prod)}

Some frameworks change scheduling and checks based on \texttt{process\.env\.NODE\_ENV}.
The benchmark Rollup config sets this value at bundle time.

\begin{itemize}
\item Rollup config: \texttt{rollup.benchmarks.config.mjs}
\item It uses \texttt{@rollup/plugin-replace} to define \texttt{process.env.NODE\_ENV}.
\end{itemize}

NPM scripts in \texttt{package.json} expose common modes:

\begin{itemize}
\item \texttt{build:bench:dev}
\item \texttt{build:bench:prod}
\end{itemize}

To be fair, comparisons should run with comparable modes (usually prod vs prod).

\subsection{What the Rollup configs actually do}

The benchmark bundles are built through two rollup configs:

\begin{itemize}
\item \texttt{rollup.benchmarks.config.mjs} (full UI runner)
\item \texttt{rollup.benchmarks.smoke.config.mjs} (perf-smoke gate)
\end{itemize}

Both use \texttt{@rollup/plugin-replace} to inject:\\
	exttt{process.env.NODE\_ENV = 'production'} by default.

This matters for fairness because some libraries (notably React) perform extra checks and different scheduling policies in development builds.

\section{Perf smoke (CI-friendly)}

For CI and ``obvious regression'' detection, the repo includes a reduced entry:

\begin{itemize}
\item \texttt{benchmarks/perf-smoke.ts}
\item \texttt{benchmarks/perf-smoke.html}
\end{itemize}

It builds via \texttt{rollup.benchmarks.smoke.config.mjs} into \texttt{benchmarks/dist-smoke/}.

\subsection{What perf-smoke is actually specifying}

You can read \texttt{benchmarks/perf-smoke.ts} as an executable specification of the harness:

\begin{itemize}
\item it obtains a profile via \texttt{getBenchmarkProfile()}
\item it forces short defaults unless overridden (default profile becomes \texttt{quick})
\item it runs a small fixed set of cases (currently \texttt{text-1m} and \texttt{list-10k-1pct}, HRBR variants)
\item it emits a JSON report that includes profile, timestamp, and result rows
\item it enforces very wide thresholds to catch catastrophic regressions
\end{itemize}

Even though this isn't a Vitest file, it plays the same role as tests in earlier chapters: it defines what "performance sanity" means for this repository.

\subsection{Why the thresholds are wide}

\texttt{perf-smoke.ts} prints a JSON report to \texttt{console.log} and then checks very wide thresholds:

\begin{itemize}
\item fail only when \texttt{avgMs > 250} or \texttt{p95Ms > 500}
\end{itemize}

This is intentional: CI machines are noisy.
The smoke test is designed to catch catastrophic regressions (infinite loops, accidental \(\mathcal{O}(n^2)\) behavior), not micro-optimizations.
This is also why the gate uses two different lenses:

\begin{itemize}
\item \texttt{avgMs}: for sustained runaway behavior
\item \texttt{p95Ms}: for severe tail latency spikes
\end{itemize}

Tail latency matters because user experience is dominated by occasional slow frames.

\section{Instrumentation and benchmarks: keep them separate}

Optional runtime instrumentation exists in \texttt{runtime/devtools.ts} (Chapter~\ref{ch:devtools}).

Methodology rule: keep instrumentation disabled during benchmarks unless you are explicitly measuring it.
Otherwise, the measurement includes the measurement.

In practice, if you want to correlate benchmark results with DOM op counts:

\begin{itemize}
\item run once with instrumentation off to measure baseline
\item run again with instrumentation enabled (and accept the overhead) to measure the profile of operations
\end{itemize}

Comparing these two runs helps you understand whether a regression is dominated by DOM work, scheduling, or allocation overhead.

\section{Checklist: how to run a fair comparison}

A practical checklist grounded in this repo's structure:

\begin{enumerate}
\item Build benchmark bundle in a known mode (dev or prod) and record it.
\item Use a named profile (quick/default/stress) so warmup and frames are consistent.
\item Fix a seed and include it in the report.
\item Run multiple times and compare distributions (p95 and drops), not just averages.
\item If a result looks surprising, isolate the case and profile it (browser performance tools).
\end{enumerate}

\section{Online resources}

\begin{itemize}
\item Measuring performance in the browser and why alternatives to \texttt{Date.now()} matter:\\
\href{https://developer.mozilla.org/en-US/docs/Web/API/Performance/now}{MDN: \texttt{performance.now()}}.

\item About frame budgets and why 16.7ms is the common baseline:\\
\href{https://web.dev/articles/rendering-performance}{web.dev: rendering performance}.

\item Why warmup matters (JIT, inline caches, tiered compilation):\\
\href{https://v8.dev/blog}{V8 blog (background on optimization and warmup effects)}.

\item Percentiles and tail latency as a reporting best practice:\\
\href{https://en.wikipedia.org/wiki/Percentile}{Wikipedia: percentile},\\
\href{https://research.google/pubs/pub40801/}{Google SRE: tail latency (The Tail at Scale)}.

\item General benchmark hygiene and common pitfalls:\\
\href{https://benchmarking.sites/}{benchmarking.sites (community notes)}.
\end{itemize}
